\chapter*{Remerciements}

Je tiens d'abord à remercier mon directeur de recherche et mentor, le professeur Michel Côté, sans qui ce travail n'existerait tout simplement pas. Michel, tu m'as permis de réaliser mon rêve d'enfance : celui de devenir chercheur scientifique. Pour cela, pour la liberté intellectuelle que tu m'as accordée, pour le soutien constant, et pour les discussions stimulantes autant sur la physique que sur la vie, je te serai éternellement reconnaissant.

Dans la même veine, je remercie tous les membres du groupe Côté, anciens comme nouveaux. En particulier, merci à Tony Vuu, Emmanuel Bourret, Surya Timilsina, Azin Kazemi Moridani, Dounia Shaaban Kabakibo et aux résidents temporaires de notre bureau, pour avoir créé un environnement respectueux, propice à la collaboration et où il fait bon vivre. On ne réussit bien que là où l'on se sent bien et je souhaite à chaque étudiant de trouver un groupe semblable.

Hors du groupe, je remercie le personnel de l'Université de Montréal, en particulier celui du campus MIL et du CEPSUM, qui m'a offert un cadre d'études et d'entraînement sportif exceptionnel. Un merci tout particulier à Sophie Tremblay du Département de physique dont l'efficacité m'a épargné bien des tracas administratifs. Merci aussi aux étudiants dont j'ai été l'auxiliaire d'enseignement : votre curiosité et votre enthousiasme m'ont fait redécouvrir ma passion pour la pédagogie.

Ce travail a bénéficié du soutien financier du Conseil de recherches en sciences naturelles et en génie du Canada (CRSNG), par une bourse du programme FONCER de l'Académie en matériaux avancés (AMA), ainsi que de subventions du CRSNG et du Fonds de recherche du Québec -- Nature et technologies (FRQNT) octroyées au professeur Côté. Les calculs numériques ont été effectués grâce aux ressources fournies par Calcul Québec et l'Alliance de recherche numérique du Canada.

En dehors du milieu universitaire, je remercie mes amis, qui se reconnaîtront, pour les aventures épiques et les soirées mémorables qui ont ponctué ces années. J'accorde une mention spéciale à celles et ceux qui m'ont accompagné dans ma pratique sportive, à l'entraînement comme au bord du parcours. Les études supérieures en physique sont un ultramarathon dont personne ne connaît ni la distance ni le dénivelé et vous avez été mes ravitaillements.

Mes derniers remerciements, et non les moindres, vont à ma famille. D'abord, merci à mes parents, Laurence et Stéphane : sans vous non plus rien de tout ça n'aurait été possible. D'une part parce que vous m'avez mis au monde et, d'autre part, parce que vous m'avez toujours encouragé à poursuivre mes rêves. C'est le plus beau cadeau que des parents puissent faire à un enfant. Enfin, je remercie mes deux frères, Louis et Raphaël, pour les discussions tout aussi variées qu'enrichissantes. Vous me comprenez mieux que quiconque et je suis profondément choyé de vous avoir dans ma vie. C'est sans doute le deuxième plus beau cadeau que des parents puissent offrir.

\bigskip

